\begin{abstract}
Recent 3D Gaussian Splatting (3DGS) Dropout methods address overfitting under sparse-view conditions by randomly nullifying Gaussian opacities.
%
However, we identify a neighbor compensation effect in these approaches: dropped Gaussians are often compensated by their neighbors, weakening the intended regularization. Moreover, these methods overlook the contribution of high-degree spherical harmonic coefficients (SH) to overfitting.
%
To address these issues, we propose \methodname{}, a novel anchor-based Dropout strategy. Rather than dropping Gaussians independently, our method randomly selects certain Gaussians as anchors and simultaneously removes their spatial neighbors. This effectively disrupts local redundancies near anchors and encourages the model to learn more robust, globally informed representations.
%
Furthermore, we extend the Dropout to color attributes by randomly dropping higher-degree SH to concentrate appearance information in lower-degree SH. This strategy further mitigates overfitting and enables flexible post-training model compression via SH truncation.
%
Experimental results demonstrate that \methodname{} substantially outperforms existing Dropout methods with negligible computational overhead, and can be readily integrated into various 3DGS variants to enhance their performances.
\end{abstract}